\documentclass[a4paper,12pt]{article}
\usepackage[utf8]{inputenc}
\usepackage[english, russian]{babel}
\usepackage{amsmath, amssymb, amsthm}
\usepackage{enumitem}%оформление списков
\usepackage{varwidth}
\usepackage{graphicx}
\usepackage{float}
\usepackage{indentfirst}
\usepackage{url}
\usepackage{seqsplit}
\usepackage{pdflscape}
\usepackage{tabularx}
\usepackage{xcolor}
\usepackage{listings}
\usepackage{mathtools}
\usepackage{hyperref}
\usepackage[backend=biber, style=numeric, sorting=none]{biblatex}
\usepackage{titlesec}
\usepackage{appendix}
\usepackage{alltt}
\usepackage[hidelinks]{hyperref}
\usepackage{booktabs}

\addbibresource{my_libs_Zaikov_I_V.bib}


\usepackage[a4paper, top=20mm, left=30mm, right=20mm, bottom=20mm]{geometry}
\newcommand{\dim}{\mathrm{dim}}
\newcommand{\lin}{\mathrm{lin}}
\newcommand{\diag}{\mathrm{diag}}
\newcommand{\R}{\mathbb{R}}
\newcommand{\Rn}{\mathbb{R^n}}

\lstset{
    language=C++,
    basicstyle=\ttfamily\small,
    keywordstyle=\color{blue},
    commentstyle=\color{green!60!black},
    stringstyle=\color{red},
    numbers=left,
    numberstyle=\tiny\color{gray},
    stepnumber=1,
    breaklines=true,
    breakatwhitespace=true,
    tabsize=4,
    showspaces=false,
    showstringspaces=false
}

\newtheoremstyle{indented}% отступ
{3pt}% над
{3pt}% под
{\itshape}% шрифт тела, переключатель
{\parindent}% отступ
{\bfseries}% шрифт заголовка
{.}% точка после заголовка
{.5em}% отступ от заголовка
{}% head spec (empty = ‘normal’)

\theoremstyle{indented}
\newtheorem{theorem}{Теорема}
\newtheorem*{theorem*}{Теорема}
\newtheorem{definition}{Определение}
\newtheorem*{definition*}{Определение}
\newtheorem{following}{Следствие}
\newtheorem*{following*}{Следствие}
\newtheorem{statement}{Утверждение}
\newtheorem*{statement*}{Утверждение}

\renewcommand{\thesection}{\arabic{section}.}
\renewcommand{\thesubsection}{\arabic{section}.\arabic{subsection}.}
\renewcommand{\thesubsubsection}{\arabic{section}.\arabic{subsection}.\arabic{subsubsection}.}

\setcounter{section}{0}
\setcounter{subsection}{0}
\setcounter{equation}{0}

\begin{document}
\pagenumbering{arabic}
\setcounter{page}{1}
\pagestyle{empty}
\begin{center}
\textbf{Министерство науки и высшего образования Российской Федерации}
\\
\vspace{0.5cm}
\textbf{Федеральное государственное бюджетное образовательное\\
учреждение высшего образования} \\
\textbf{«Ярославский государственный университет им. П.Г. Демидова»}

\vspace{0.5cm}

Кафедра математического анализа
\vfill

\begin{flushright}
  Сдано на кафедру\\
  « \underline{\hspace{0.8cm}} » \underline{\hspace{3cm}} 2026 г.\\
  Заведующий кафедрой \\
  \underline{д. ф.-м. н., доцент }\\
  \underline{\hspace{3.4cm}} Невский М.В.\\
\end{flushright}
\vfill

Выпускная квалификационная работа магистра \\

\vspace{0.5cm}

\textbf{Задача о минимальной норме интерполяционного проектора} \\
(Направление подготовки магистров 01.03.02 Прикладная математика и информатика)
\vfill


\bigskip


\end{center}

\begin{flushright}
Научный руководитель\\
\underline{к. ф.-м. н., доцент }\\
\underline{\hspace{3.4cm}} Ухалов А.Ю.\\
« \underline{\hspace{0.8cm}} » \underline{\hspace{3cm}} 2026 г.\\
\vspace{1.5cm}
Студент группы \underline{ПМИ-21МО}\\
\underline{\hspace{3.5cm}} Зайков И.В. \\
« \underline{\hspace{0.8cm}} » \underline{\hspace{3cm}} 2026 г.\\
\end{flushright}

\vspace{0.5cm}
\begin{center}
Ярославль 2026 г.
\end{center}
\newpage

\section*{Реферат}
\addcontentsline{section}{Реферат}

Объем 51 страница, 5 параграфов, 2 таблицы, 3 приложения. Список цитированной литературы содержит 7 источников.
\par
Изучается интерполяция непрерывной функции $f(x)$, определенной на кубе
$Q_n = [0, 1]^n$ , с помощью полиномов от $n$ переменных степени не выше 1. Ставится
задача оценивания минимальной нормы проектора для $n = 31$. Верхняя оценка $\theta_{31}$
получается из рассмотрения проектора с узлами интерполяции, совпадающими с
вершинами симплекса максимального объема в $Q_n$, построенного на основе матрицы Адамара. Получена новая оценка $\theta_{31} \leq 5.0$.
\par
\textbf{Ключевые слова: интерполяция, интерполяционный проектор, норма, матрица Адамара, симплекс максимального объема.}

\newpage

\pagestyle{plain}
\tableofcontents

\newpage

% ============================================================================
\section*{Введение}
\addcontentsline{toc}{section}{Введение}

В выпускной квалификационной работе рассматривается задача о минимальной норме интерполяционного проектора. Изучается интерполяция непрерывной функции $f(x)$, определённой на кубе $Q_n = [0,1]^n$, с помощью полиномов от $n$ переменных степени не выше $1$.
\par
Для построения интерполяционного полинома необходимо задать $n+1$ узел — вершины невырожденного симплекса $S \subset Q_n$.
\par
Интерполяционный проектор $P$ ставит в соответствие функции $f$ полином $P f$, совпадающий с $f$ в узлах. Норма оператора $P$ характеризует устойчивость интерполяции к погрешностям входных данных: чем меньше норма, тем ближе интерполяционный полином к наилучшему приближению.

Минимальная норма интерполяционного проектора $\theta_n$ определяется как наименьшее возможное значение $\|P\|$ при условии, что узлы интерполяции принадлежат кубу $Q_n$ и образуют невырожденный симплекс. Задача нахождения $\theta_n$ является классической в теории приближений и имеет важное практическое значение, поскольку позволяет выбрать расположение узлов, минимизирующее множитель $(1 + \|P\|)$ в неравенстве Лебега и, следовательно, дающее наилучшую гарантированную оценку погрешности интерполяции.

В настоящей работе для получения верхней оценки $\theta_n$ используется подход, основанный на рассмотрении симплексов максимального объёма, вписанных в гиперкуб $Q_n$. Такие симплексы могут быть построены из матриц $(-1/1)$ порядка $n+1$, определители которых максимальны, то есть из матриц Адамара. Норма проектора для симплекса вычисляется по формуле
\[
\|P\| = \max_{x \in \operatorname{ver}(Q_n)} \sum_{j=1}^{n+1} |\lambda_j(x)|,
\]
где $\lambda_j$ — базисные многочлены Лагранжа (барицентрические координаты). Для нахождения максимума требуется перебрать все $2^n$ вершин гиперкуба.
\par
Для размерности $n = 31$ было исследовано шесть типов матриц Адамара порядка $32$ (\texttt{syl}, \texttt{pal}, \texttt{t1}, \texttt{t2}, \texttt{t3}, \texttt{t4}). Для каждого типа построен симплекс максимального объёма $S \subset Q_{31}$ и вычислена норма интерполяционного проектора $\|P_S\|$. В качестве оценки сверху для $\theta_{31}$ принимается минимум норм по всем рассмотренным симплексам:
\[
\theta_{31} \le \min_{S} \|P_S\|.
\]
\par
Была разработана программа на языке C++ с использованием библиотеки Eigen для линейно-алгебраических операций. Для ускорения вычислений использовались параллельные вычисления на $16$ потоках, что позволило завершить полный перебор за $290$ секунд (около $5$ минут).
\par
Для верификации программной реализации был проведён ряд проверок:
\begin{itemize}
    \item проверка интерполяционных свойств $\lambda_j(x^{(k)}) = \delta_{jk}$ с точностью $10^{-10}$;
    \item верификация на пространствах $\mathbb{R}^n$ меньших размерностей путём сравнения с известными из литературы результатами ($n = 0, 1, 3, 7, 11, 15, 19, 23, 27$);
    \item повторная проверка найденного максимума после завершения перебора.
\end{itemize}

В результате работы получена верхняя оценка минимальной нормы интерполяционного проектора для размерности $n = 31$:
\[
\theta_{31} \le 5.000.
\]

\newpage

% ============================================================================
\section{Используемые обозначения}

Везде далее предполагается, что $n$ — натуральное число ($n \in \mathbb{N}$). Обозначим через $Q_n$ единичный куб в $\mathbb{R}^n$:
\[
Q_n = [0,1]^n = \{(x_1, \dots, x_n) : 0 \le x_i \le 1,\ i = 1,\dots,n\}.
\]

Через $C(Q_n)$ обозначим пространство непрерывных функций $f : Q_n \to \mathbb{R}$ с равномерной нормой
\[
\|f\|_{C(Q_n)} = \max_{x \in Q_n} |f(x)|.
\]

Под $\Pi_1(\mathbb{R}^n)$ понимается совокупность многочленов от $n$ переменных степени $\le 1$, т.е. линейная оболочка
\[
\Pi_1(\mathbb{R}^n) = \operatorname{span}\{1, x_1, x_2, \dots, x_n\}.
\]

Симплекс $S \subset \mathbb{R}^n$ с вершинами $x^{(1)}, \dots, x^{(n+1)}$ обозначается
\[
S = \operatorname{conv}\{x^{(1)}, \dots, x^{(n+1)}\}.
\]

% ============================================================================
\section{Интерполяционный проектор}

\subsection{Постановка задачи интерполяции}

В \cite{nevsky_geo} описывается задача линейной интерполяции на $n$-мерном кубе.
\par
Пусть $Q_n = [0,1]^n$ — единичный куб в $\mathbb{R}^n$, а $C(Q_n)$ — пространство непрерывных функций $f: Q_n \to \mathbb{R}$ с равномерной нормой
\[
\|f\|_{C(Q_n)} = \max_{x \in Q_n} |f(x)|.
\]

Рассмотрим интерполяцию функций из $C(Q_n)$ полиномами от $n$ переменных степени не выше $1$. Пространство таких полиномов обозначим через $\Pi_1(\mathbb{R}^n)$:
\[
\Pi_1(\mathbb{R}^n) = \operatorname{span}\{1, x_1, x_2, \dots, x_n\}.
\]

Для построения интерполяционного полинома необходимо задать $n+1$ узел
\[
x^{(1)}, x^{(2)}, \dots, x^{(n+1)} \in Q_n,
\]
образующих невырожденный симплекс $S$. Требуется найти полином $p \in \Pi_1(\mathbb{R}^n)$ такой, что
\[
p(x^{(j)}) = f(x^{(j)}), \quad j = 1, \dots, n+1.
\]

\subsection{Базисные многочлены Лагранжа и барицентрические координаты}

В \cite{nevsky_geo} описан способ построения базисных многочленов Лагранжа с помощью барицентрических координат, определённых относительно вершин симплекса.

Обозначим вершины симплекса $S$ через
\[
x^{(j)} = \left(x_1^{(j)}, x_2^{(j)}, \dots, x_n^{(j)}\right), \quad 1 \le j \le n+1.
\]

Рассмотрим матрицу $A$ размера $(n+1) \times (n+1)$, в которой \textbf{строки} соответствуют вершинам симплекса, дополненным единицей в последнем столбце:
\[
A = \begin{pmatrix}
x_1^{(1)} & x_2^{(1)} & \dots & x_n^{(1)} & 1 \\
x_1^{(2)} & x_2^{(2)} & \dots & x_n^{(2)} & 1 \\
\vdots & \vdots & & \vdots & \vdots \\
x_1^{(n+1)} & x_2^{(n+1)} & \dots & x_n^{(n+1)} & 1
\end{pmatrix}.
\]

Столбец единиц введён для того, чтобы учесть условие нормировки барицентрических координат $\sum_{j=1}^{n+1} \lambda_j(x) = 1$. Переход к однородным координатам позволяет записать систему для нахождения барицентрических координат в матричной форме.
\par
Рассмотрим матрицу $A$ размера $(n+1) \times (n+1)$, строки которой соответствуют вершинам симплекса, дополненным единицей в последнем столбце:
\[
A = \begin{pmatrix}
x_1^{(1)} & x_2^{(1)} & \dots & x_n^{(1)} & 1 \\
x_1^{(2)} & x_2^{(2)} & \dots & x_n^{(2)} & 1 \\
\vdots & \vdots & & \vdots & \vdots \\
x_1^{(n+1)} & x_2^{(n+1)} & \dots & x_n^{(n+1)} & 1
\end{pmatrix}.
\]
\par
Тогда для точки $x = (x_1, \dots, x_n)$ возможно найти ее барицентрические координаты $\lambda(x) = (\lambda_1(x), \dots, \lambda_{n+1}(x))^T$ из системы
\[
A \lambda(x) = b(x),
\]
где $b(x) = (x_1, \dots, x_n, 1)^T$ — расширенный вектор точки. Поскольку симплекс невырожден, матрица $A$ обратима, и решение системы имеет вид
\[
\lambda(x) = A^{-1} b(x).
\]
\par
Пусть $\Delta$ - определитель матрицы $A$. Рассмотрим $\Delta_j(x)$ определитель, который получается из $\Delta$ заменой $j$-й строки на строку $(x_1, x_2, \dots, x_n, 1)$. Тогда функции
\[
\lambda_j(x) = \frac{\Delta_j(x)}{\Delta}, \quad j = 1, \dots, n+1,
\]
являются также барицентрическими координатами точки $x$ относительно симплекса $S$. Они обладают следующими свойствами:
\begin{enumerate}
    \item $\lambda_j(x) \in \Pi_1(\mathbb{R}^n)$, то есть являются линейными функциями;
    \item $\lambda_j(x^{(k)}) = \delta_{jk}$, где $\delta_{jk}$ — символ Кронекера;
    \item $\sum_{j=1}^{n+1} \lambda_j(x) = 1$ для всех $x \in \mathbb{R}^n$.
\end{enumerate}

Таким образом, барицентрические координаты $\lambda_j(x)$ являются \textbf{базисными многочленами Лагранжа}, соответствующими симплексу $S$. Коэффициенты этих многочленов составляют $j$-й столбец обратной матрицы $A^{-1}$. Действительно, пусть $A^{-1} = (l_{ij})_{i,j=1}^{n+1}$. Тогда
\[
\lambda_j(x) = l_{1j} x_1 + l_{2j} x_2 + \dots + l_{nj} x_n + l_{n+1,j}.
\]

Любой полином $p \in \Pi_1(\mathbb{R}^n)$ может быть разложен по этому базису:
\[
p(x) = \sum_{j=1}^{n+1} p(x^{(j)}) \lambda_j(x). \tag{1}
\]

\subsection{Интерполяционный проектор}

Для заданной функции $f \in C(Q_n)$ определим отображение $P: C(Q_n) \to \Pi_1(\mathbb{R}^n)$ формулой
\[
P f(x) = \sum_{j=1}^{n+1} f(x^{(j)}) \lambda_j(x). \tag{2}
\]

Из свойства $\lambda_j(x^{(k)}) = \delta_{jk}$ следует, что
\[
P f(x^{(j)}) = f(x^{(j)}), \quad j = 1, \dots, n+1,
\]
то есть полином $P f$ интерполирует функцию $f$ в узлах $x^{(j)}$. Отображение $P$ называется \textbf{интерполяционным проектором}.

Оператор $P$ линеен и обладает свойством идемпотентности: для любого полинома $p \in \Pi_1(\mathbb{R}^n)$ выполнено $P p = p$.

\subsection{Норма интерполяционного проектора}

Норма оператора $P$ определяется стандартным образом:
\[
\|P\| = \sup_{\|f\|_{C(Q_n)} \le 1} \|P f\|_{C(Q_n)}.
\]

В \cite{nevsky_geo} показано, что для гиперкуба $Q_n$ максимум суммы модулей базисных функций достигается в одной из его вершин, т.е. справедлива формула:
\[
\|P\| = \max_{x \in \operatorname{ver}(Q_n)} \sum_{j=1}^{n+1} |\lambda_j(x)|. \tag{3}
\]

\subsection{Неравенство Лебега}

% todo: ссылку на литературу
Пусть $E_n(f)$ — наилучшее приближение функции $f$ полиномами из $\Pi_1(\mathbb{R}^n)$:
\[
E_n(f) = \min_{p \in \Pi_1(\mathbb{R}^n)} \|f - p\|_{C(Q_n)}.
\]

Для интерполяционного проектора $P$ справедливо \textbf{неравенство Лебега}:
\[
\|f - P f\|_{C(Q_n)} \le (1 + \|P\|) \cdot E_n(f). \tag{4}
\]

Это неравенство показывает, что погрешность интерполяции не превосходит наилучшего приближения, умноженного на множитель $(1 + \|P\|)$. Таким образом, чем меньше норма проектора, тем ближе интерполяционный полином к наилучшему приближению.

\subsection{Минимальная норма проектора}

Обозначим через $\theta_n$ минимальную норму интерполяционного проектора при условии, что узлы интерполяции принадлежат кубу $Q_n$ и образуют невырожденный симплекс:
\[
\theta_n = \min_{\substack{x^{(j)} \in Q_n \\ \det(A) \neq 0}} \|P\|.
\]
\par
Задача о нахождении $\theta_n$ является классической в теории приближений. Она сводится к поиску такого набора узлов (симплекса $S \subset Q_n$), который минимизирует правую часть (3). Отметим, что выражение под знаком минимума зависит от $n(n+1)$ переменных (координат $n+1$ вершины симплекса), а вычисление $\lambda_j(x)$ требует обращения матрицы $A$, что делает задачу крайне трудоёмкой для больших $n$.
\par
Задача настоящей работы состоит в получении верхней оценки для $\theta_{31}$. Для этого рассматривается конечное множество симплексов $\mathcal{S}$, построенных по шести матрицам Адамара порядка $32$. Для каждого $S \in \mathcal{S}$ вычисляется норма $\|P_S\|$, после чего выбирается минимальное значение:
\[
\theta_{31} \le \min_{S \in \mathcal{S}} \|P_S\|
\]

\subsection{Оценки нормы проектора}

Для величины $\theta_n$ известны следующие общие оценки (см. \cite{nevsky_geo}):
\[
\frac{1}{4}\sqrt{n} < \theta_n < 3\sqrt{n}, \qquad 3 - \frac{4}{n+1} \le \theta_n.
\]

Если $S$ — симплекс максимального объёма в $Q_n$, а $P$ — интерполяционный проектор с узлами в вершинах $S$, то справедливо асимптотическое равенство
\[
\|P\| \asymp \theta_n,
\]
где запись $f(n) \asymp g(n)$ означает, что существуют константы $c_1, c_2 > 0$, не зависящие от $n$, такие что
\[
c_1 g(n) \le f(n) \le c_2 g(n) \quad \text{для всех достаточно больших } n.
\]

Таким образом, норма проектора по узлам в вершинах симплекса максимального объёма и величина $\theta_n$ имеют одинаковый порядок роста. Это свойство лежит в основе метода уточнения оценок для $\theta_n$, применяемого в настоящей работе. В качестве верхней границы $\theta_n$ можно взять значение нормы проектора с узлами в вершинах симплекса максимального объёма в $Q_n$.
\par
В \cite{nevsky_ukhalov_sel_problems} приведены оценки $\theta_n$ для различных размерностей. Наиболее точные известные оценки для $1 \le n \le 26$ систематизированы в таблице ~\ref{tab:theta_estimates}.
% ============================================================================
\section{Матрицы Адамара}

\subsection{Определение и основные свойства}

\textbf{($-1/1$)-матрицей} называется матрица, каждый элемент которой равен $-1$ или $1$. \textbf{Матрицей Адамара} порядка $n$ называется квадратная $(-1/1)$-матрица $H$ размера $n \times n$, удовлетворяющая условию:
\[
H H^T = n I_n,
\]
где $I_n$ — единичная матрица. Из этого свойства непосредственно следует, что строки (и столбцы) матрицы Адамара попарно ортогональны. Для определителя матрицы Адамара справедливо соотношение:
\[
|\det H| = n^{n/2}.
\]

Подробная информация о матрицах Адамара содержится в книге \cite{hall_comb_1970}

Матрицы Адамара существуют не для всех порядков. Необходимым условием существования является $n = 1$, $n = 2$ или $n$, кратное $4$. Достаточность этого условия составляет содержание \textbf{гипотезы Адамара}, которая до настоящего времени не доказана. Тем не менее, для огромного количества порядков, кратных $4$, матрицы Адамара построены. В частности, для порядка $32 = 4 \cdot 8$ матрицы Адамара существуют.

\subsection{Количество неэквивалентных матриц Адамара}

Две матрицы Адамара называются \textbf{эквивалентными}, если одна может быть получена из другой с помощью перестановок строк и столбцов и умножения строк и столбцов на $-1$. Количество неэквивалентных матриц Адамара для малых порядков приведено в таблице~\ref{tab:hadamard_counts}.

% todo: сформулировать
Из таблицы ~\ref{tab:hadamard_counts} видно, что количество матриц быстро растёт с увеличением порядка. Для $n = 32$ оно составляет уже более тринадцати миллионов, что делает полный перебор всех симплексов максимального объёма в $\mathbb{R}^{31}$ практически невозможным.

\subsection{Методы построения матриц Адамара}

Существует несколько различных методов построения матриц Адамара. В настоящей работе используются матрицы порядка $32$, взятые из открытой библиотеки \cite{sloane_hadamard}. Перечень использованных типов матриц приведён ниже:

\subsubsection{Метод Сильвестра (тип `syl`)}

Для порядков, являющихся степенью двойки, матрицы Адамара могут быть построены рекурсивным методом Сильвестра. Рекуррентное соотношение имеет вид:
\[
H_{1} = (1), \qquad H_{2n} = \begin{pmatrix}
H_n & H_n \\
H_n & -H_n
\end{pmatrix}.
\]

Например, для $n = 2$ получаем:
\[
H_2 = \begin{pmatrix}
1 & 1 \\
1 & -1
\end{pmatrix},
\]
для $n = 4$:
\[
H_4 = \begin{pmatrix}
1 & 1 & 1 & 1 \\
1 & -1 & 1 & -1 \\
1 & 1 & -1 & -1 \\
1 & -1 & -1 & 1
\end{pmatrix}.
\]

Матрицы, построенные этим методом, обозначаются суффиксом \texttt{syl} (например, \texttt{had.32.syl}).

\subsubsection{Метод Пэли (тип `pal`)}

Конструкция Пэли (Paley construction) позволяет строить матрицы Адамара для порядков $n = q + 1$, где $q$ — степень нечётного простого числа. В основе метода лежат свойства квадратичных вычетов в конечных полях. Для порядка $32 = 31 + 1$ существует матрица Адамара типа Пэли, обозначаемая суффиксом \texttt{pal} (например, \texttt{had.32.pal}).

\subsubsection{Другие типы матриц порядка 32}

Для порядка $32$ кроме матриц типов \texttt{syl} и \texttt{pal} существуют также матрицы, обозначаемые \texttt{t1}, \texttt{t2}, \texttt{t3}, \texttt{t4} (например, \texttt{had.32.t1}, \texttt{had.32.t2}, \texttt{had.32.t3}, \texttt{had.32.t4}). Эти матрицы получены методом тензорного произведения. В настоящей работе исследуются все перечисленные типы матриц для сравнения значений норм интерполяционных проекторов, получаемых на соответствующих симплексах.

\subsection{Переход от $(-1/1)$-матриц к $(0/1)$-матрицам максимального определителя}

Как показано в \cite{nevsky_geo}, симплекс максимального объёма в гиперкубе $Q_n = [0,1]^n$ может быть получен из $(0/1)$-матрицы порядка $n$, имеющей максимальный определитель. В качестве первых $n$ вершин такого симплекса можно взять строки этой матрицы, а в качестве последней $(n+1)$-й вершины добавить нулевую строку.
Максимальный $(0/1)$-определитель порядка $n$ может быть получен из максимального $(-1/1)$-определителя порядка $n+1$ с помощью следующей процедуры. Пусть $A$ — максимальный $(-1/1)$-определитель порядка $n+1$.

\begin{enumerate}
    \item Каждый столбец $A$, начинающийся с $-1$, умножается на $-1$.
    \item Каждая строка новой матрицы, начинающаяся с $-1$, также умножается на $-1$.
\end{enumerate}

Матрицу порядка $n+1$, которая получится в результате выполнения шагов 1–2, обозначим через $B$. У этой $(-1/1)$-матрицы первый столбец и первая строка целиком состоят из $1$.

\begin{enumerate}
    \setcounter{enumi}{2}
    \item Пусть $C$ — подматрица порядка $n$ матрицы $B$, стоящая в строках и столбцах с номерами $2, \dots, n+1$. В этой матрице производится следующая замена элементов: $1$ заменяется на $0$, а $-1$ на $1$. Получившуюся $(0/1)$-матрицу порядка $n$ обозначим $D$.
\end{enumerate}

Полученный таким образом определитель $D$ является максимальным $(0/1)$-определителем порядка $n$.

\subsection{Построение симплекса максимального объёма в $\mathbb{R}^{31}$}

В настоящей работе для размерности $n = 31$ в качестве исходных берутся матрицы Адамара порядка $32$ следующих типов: \texttt{syl}, \texttt{pal}, \texttt{t1}, \texttt{t2}, \texttt{t3}, \texttt{t4}. Для каждой из них выполняется описанная выше процедура перехода к $(0/1)$-матрице $D$ порядка $31$, после чего строки $D$ становятся первыми $31$ вершиной симплекса, а в качестве последней, $32$-й вершины добавляется нулевая строка:
\[
V = \{0, \operatorname{row}_1(D), \operatorname{row}_2(D), \dots, \operatorname{row}_{31}(D)\} \subset [0,1]^{31}.
\]

Таким образом, для каждого типа матрицы Адамара получается соответствующий симплекс максимального объёма, вписанный в гиперкуб $[0,1]^{31}$. Все вершины таких симплексов являются вершинами гиперкуба. Полученные симплексы используются в настоящей работе для вычисления верхних оценок минимальной нормы интерполяционного проектора $\theta_{31}$.

\subsection{Замечание о количестве симплексов}

Для размерности $n = 31$ количество неэквивалентных матриц Адамара порядка $32$ составляет $13\,710\,027$ \cite{kharaghani2013hadamard}. Это означает, что существует такое же количество различных симплексов максимального объёма, вписанных в гиперкуб $Q_{31}$. Значения норм интерполяционных проекторов для разных симплексов могут различаться.
\par
В настоящей работе для исследования используются шесть матриц Адамара порядка $32$, доступных на сайте библиотеки Нила Слоуна \cite{sloane_hadamard}:
\[
\texttt{had.32.syl},\quad \texttt{had.32.pal},\quad \texttt{had.32.t1},\quad \texttt{had.32.t2},\quad \texttt{had.32.t3},\quad \texttt{had.32.t4}.
\]

Полный перебор всех $13\,710\,027$ симплексов является крайне трудоёмкой задачей и не входит в рамки настоящей работы. Исследование ограничивается указанными шестью представителями, что позволяет получить содержательные оценки при разумных вычислительных затратах.
% ============================================================================
\section{Описание программной реализации}

Исходный код программы доступен в открытом репозитории по адресу:
\url{https://gitflic.ru/project/zaykov_ivan/minimal_norm_calculator.git} (или \url{https://github.com/Ivan-Zaykov/minimal_norm_calculator.git}).

Программная реализация выполнена на языке C++ (стандарт C++17). Для линейно-алгебраических операций используется библиотека Eigen \cite{eigen_lib}. Архитектура программы позволяет легко адаптировать её для вычисления нормы интерполяционного проектора для любых симплексов максимального объёма, не только построенных из матриц Адамара. Сборка проекта осуществляется с помощью CMake.

\subsection{Общая архитектура}

Программа построена в соответствии с принципами SOLID. Основные классы:

\begin{itemize}
    \item \texttt{ISimplexProvider} — интерфейс для получения вершин симплекса;
    \item \texttt{SilvesterMethodProvider} — построение симплекса из матрицы Адамара методом Сильвестра;
    \item \texttt{HadamardFileProvider} — построение симплекса по матрице Адамара из файла;
    \item \texttt{SimplexProviderFactory} — фабрика для создания объектов классов, отвечающих за построение симплекса;
    \item \texttt{HadamardMatrixIterator} — итератор для последовательного чтения матриц из файла;
    \item \texttt{LebesgueFunction} — вычисление значения функции Лебега (суммы модулей базисных многочленов) в точке;
    \item \texttt{FullEnumUpperBoundCalculator} — многопоточный перебор вершин гиперкуба;
    \item \texttt{CommandLineParser} — разбор аргументов командной строки;
    \item \texttt{Processing} — организация основных вычислительных режимов.
\end{itemize}

\subsection{Ввод и вывод данных}

Программа поддерживает два режима работы:
\begin{enumerate}
    \item Построение симплекса по матрице Адамара Сильвестра (по умолчанию);
    \item Загрузка матриц Адамара из файла (символьный или числовой формат);
\end{enumerate}

Файлы с матрицами Адамара могут быть загружены с сайта библиотеки Нила Слоуна \cite{sloane_hadamard}. Для автоматической загрузки всех доступных матриц порядка 32 разработан скрипт \texttt{download\_hadamard.sh}. Он скачивает файлы \texttt{had.32.syl}, \texttt{had.32.pal}, \texttt{had.32.t1} – \texttt{had.32.t4} и объединяет их в один файл с пустыми строками-разделителями. Кроме того, в репозитории программы в каталоге \texttt{preset} уже содержатся готовые матрицы различных размерностей, что позволяет использовать их без дополнительной загрузки.

Управление режимами работы осуществляется через аргументы командной строки:
\begin{itemize}
    \item \texttt{--dimension N} — установить размерность пространства;
    \item \texttt{--load-hadamard FILE} — загрузить матрицы Адамара из файла;
    \item \texttt{--batch DIR} — пакетная обработка всех файлов в каталоге;
    \item \texttt{--threads N} — количество потоков (0 = автоматически);
    \item \texttt{--log-interval N} — интервал логирования (количество вершин);
    \item \texttt{--range START END} — частичный перебор вершин в заданном диапазоне.
\end{itemize}


\subsection{Дополнительные проверки корректности построения проектора}

Для верификации программной реализации и достоверности полученных результатов был проведён ряд дополнительных проверок.

\subsubsection{Проверка интерполяционных свойств}

Для каждого построенного симплекса выполнялась проверка условия
\[
\lambda_j(x^{(k)}) = \delta_{jk}, \qquad j,k = 1,\dots,n+1,
\]
где $\lambda_j$ — базисные многочлены Лагранжа (барицентрические координаты), а $x^{(k)}$ — вершины симплекса. Проверка проводилась с точностью $10^{-5}$. Для всех рассмотренных матриц Адамара максимальная ошибка не превысила $10^{-10}$, что подтверждает корректность построения симплексов и вычисления базисных функций.

\subsubsection{Проверка полного перебора}

Для малых размерностей ($n = 0, 1, 3, 7, 11, 15, 19, 23, 27$) был выполнен полный перебор всех вершин гиперкуба, и полученные оценки нормы проектора были сопоставлены с известными из литературы. Результаты сравнения соответствуют приведенным в ~\ref{tab:theta_estimates}.

\subsubsection{Повторная проверка найденного максимума}

После завершения полного перебора для каждой матрицы значение нормы проектора повторно вычислялось в найденной вершине. Расхождение между значением, зафиксированным во время перебора, и значением, полученным при повторной проверке, не превышало $10^{-10}$, что подтверждает корректность синхронизации глобального максимума при многопоточном переборе.

\subsection{Вычисление нормы проектора}

Норма интерполяционного проектора вычисляется по формуле (2.3):
\[
\|P\| = \max_{x \in \operatorname{ver}(Q_n)} \sum_{j=1}^{n+1} |\lambda_j(x)|.
\]

Для нахождения максимума реализован полный перебор всех $2^{31}$ вершин гиперкуба. Каждая вершина кодируется 31-битовым целым числом, что позволяет эффективно генерировать координаты с помощью битовых операций.

Вычисления выполняются в несколько потоков (количество потоков может быть задано параметром \texttt{--threads}). Для синхронизации глобального максимума используются атомарные операции. Программа выводит прогресс выполнения каждые $10^7$ вершин с оценкой времени до завершения.

По умолчанию все результаты выводятся в стандартный поток вывода (консоль). Для сохранения результатов в файл при одновременном отображении в консоли может быть использована утилита \texttt{tee}, например:
\begin{verbatim}
 ./HadamardSimplexNorm --dimension 31 | tee results.txt
\end{verbatim}
При указании опции \texttt{--output FILE} программа самостоятельно сохраняет результаты в указанный файл.

Полный исходный код программы приведён в приложении.
% ============================================================================
\section{Результаты вычисления верхней оценки нормы}

\subsection{Параметры вычислительного эксперимента}

Вычисления проводились на следующем оборудовании и программном обеспечении:

\begin{itemize}
    \item \textbf{Процессор}: 12th Gen Intel(R) Core(TM) i5-1240P (16 ядер)
    \item \textbf{Оперативная память}: 15 ГБ
    \item \textbf{Операционная система}: Ubuntu 24.04.3 LTS
    \item \textbf{Компилятор}: GCC 13.3.0
    \item \textbf{Библиотека}: Eigen 3
    \item \textbf{Флаги компиляции}: \texttt{-O3 -march=native}
    \item \textbf{Количество потоков}: 16
    \item \textbf{Интервал логирования}: $10^7$ вершин
\end{itemize}

Полный перебор всех $2^{31}$ вершин одного гиперкуба занял 290 секунд. Средняя скорость обработки составила около 7.4 миллионов вершин в секунду.

\subsection{Результаты для размерности $n = 31$}

В результате вычислительного эксперимента для шести симплексов, построенных на основе матрицы Адамара получена следующая оценка нормы проектора:
\[
\|P\| = 5.000.
\]

\subsection{Сравнение с теоретической оценкой}

Теоретическая верхняя граница для минимальной нормы интерполяционного проектора в размерности $n = 31$ составляет:
\[
\theta_{31} \le \sqrt{32} \approx 5.656.
\]

Полученное в работе значение $\theta_{31} \le 5.000$ ниже теоретической оценки. Таким образом, наши вычисления позволили уточнить существующую оценку.
% ============================================================================

\newpage

\section*{Заключение}
\addcontentsline{toc}{section}{Заключение}

В ходе выполнения выпускной квалификационной работы была разработана программа на языке C++, предназначенная для вычисления верхних оценок минимальной нормы интерполяционного проектора $\theta_n$ для симплексов, построенных из матриц Адамара. Программа поддерживает:

\begin{itemize}
    \item построение симплекса из матрицы Адамара методом Сильвестра;
    \item загрузку матриц Адамара из файлов (символьный и числовой форматы);
    \item многопоточный полный перебор всех $2^{31}$ вершин гиперкуба;
    \item последовательная обработка нескольких файлов с данными матриц Адамара;
    \item логирование прогресса с выводом текущего результата и оценкой времени до завершения;
    \item сохранение результатов в выходной файл.
\end{itemize}

\subsection*{Полученные результаты}

В результате вычислительного эксперимента для шести симплексов, построенных по матрицам Адамара порядка $32$ различных типов, получена верхняя оценка
\[
\theta_{31} \le 5.000,
\]
Теоретическая верхняя граница для данной размерности составляет $\sqrt{32} \approx 5.656$. Таким образом, наши вычисления позволили уточнить существующую оценку.

\subsection*{Трудоёмкость анализа следующей размерности}

Следующая размерность, для которой существует матрица Адамара, — $n = 35$ (матрица порядка $36$). В этом случае гиперкуб содержит $2^{35} \approx 3.44 \times 10^{10}$ вершин. При сохранении скорости обработки $7.4 \times 10^6$ вершин в секунду полный перебор занял бы около $1.3$ часа, что выполнимо на современном компьютере. Таким образом, разработанный метод полного перебора остаётся применимым и для размерности $n = 35$.

Исходный код программы доступен в открытом репозитории по адресу:
\url{https://gitflic.ru/project/zaykov_ivan/minimal_norm_calculator.git} (или \url{https://github.com/Ivan-Zaykov/minimal_norm_calculator.git}).

% ============================================================================

\newpage

\begin{table}[h]
\centering
\caption{Количество неэквивалентных матриц Адамара}
\label{tab:hadamard_counts}
\begin{tabular}{|c|c|}
\hline
\textbf{Порядок $n$} & \textbf{Количество неэквивалентных матриц} \\
\hline
1 & 1 \\
\hline
2 & 1 \\
\hline
4 & 1 \\
\hline
8 & 1 \\
\hline
12 & 1 \\
\hline
16 & 5 \\
\hline
20 & 3 \\
\hline
24 & 60 \\
\hline
28 & 487 \\
\hline
32 & 13\,710\,027 \\
\hline
\end{tabular}
\end{table}

\newpage

\begin{table}[h]
\centering
\caption{Верхние оценки минимальной нормы интерполяционного проектора $\theta_n$ для $1 \le n \le 27$}
\label{tab:theta_estimates}
\begin{tabular}{|c|c|c|}
\hline
$n$ & $\min\|P\|$ & $\theta_n \le$ \\
\hline
1 & 1 & 1 \\
\hline
2 & 3 & $\frac{2\sqrt{5}}{5}+1 = 1.8944\ldots$ \\
\hline
3 & 2 & 2 \\
\hline
4 & $\frac{7}{3} = 2.3333\ldots$ & $\frac{3(4+\sqrt{2})}{7} = 2.3203\ldots$ \\
\hline
5 & $\frac{13}{5} = 2.6$ & $2.448804$ \\
\hline
6 & 3 & $2.60004\ldots$ \\
\hline
7 & $\frac{5}{2} = 2.5$ & $\frac{5}{2} = 2.5$ \\
\hline
8 & $\frac{22}{7} = 3.1428\ldots$ & $\frac{22}{7} = 3.1428\ldots$ \\
\hline
9 & 3 & 3 \\
\hline
10 & $\frac{19}{5} = 3.8$ & $\frac{19}{5} = 3.8\ldots$ \\
\hline
11 & 3 & 3 \\
\hline
12 & $\frac{17}{5} = 3.4$ & $\frac{17}{5} = 3.4$ \\
\hline
13 & $\frac{49}{13} = 3.7692\ldots$ & $\frac{49}{13} = 3.7692\ldots$ \\
\hline
14 & $\frac{21}{5} = 4.2$ & $\frac{21}{5} = 4.2\ldots$ \\
\hline
15 & $\frac{7}{2} = 3.5$ & $\frac{7}{2} = 3.5$ \\
\hline
16 & $\frac{21}{5} = 4.2$ & $\frac{21}{5} = 4.2$ \\
\hline
17 & $\frac{139}{34} = 4.0882\ldots$ & $\frac{139}{34} = 4.0882\ldots$ \\
\hline
18 & $\frac{95}{17} = 5.5882\ldots$ & $5.14006\ldots$ \\
\hline
19 & 4 & 4 \\
\hline
20 & $\frac{137}{29} = 4.7241\ldots$ & $4.68879\ldots$ \\
\hline
21 & $\frac{251}{50} = 5.02$ & $\frac{251}{50} = 5.02$ \\
\hline
22 & $\frac{1817}{335} = 5.4238\ldots$ & $\frac{1817}{335} = 5.4238\ldots$ \\
\hline
23 & $\frac{9}{2} = 4.5$ & $\frac{9}{2} = 4.5$ \\
\hline
24 & $\frac{103}{21} = 4.9047\ldots$ & $\frac{103}{21} = 4.9047\ldots$ \\
\hline
25 & 5 & 5 \\
\hline
26 & $\frac{474}{91} = 5.2087\ldots$ & $\frac{474}{91} = 5.2087\ldots$ \\
\hline
27 & $5.0$ & $5.0$ \\
\hline
\end{tabular}
\end{table}
\newpage

\newpage

\printbibliography

\newpage

% ============================================================================

\section*{Приложение A. Структура исходного кода}
\addcontentsline{toc}{section}{Приложение A. Структура исходного кода}

Программа реализована на языке C++ (стандарт C++17). Для линейно-алгебраических операций используется библиотека Eigen. Сборка осуществляется с помощью CMake. Ниже приведена структура исходного кода:

\begin{verbatim}
├── main.cpp                     # точка входа в программу

include/
├── basis/
│   └── SimplexBasis.h           # базис Лагранжа (барицентрические координаты)
├── cli/
│   └── CommandLineParser.h      # разбор аргументов командной строки
├── hadamard_utils/
│   ├── HadamardMatrixIterator.h # итератор для чтения матриц Адамара из файла
│   └── HadamardUtils.h          # вспомогательные функции для матриц Адамара
├── processing/
│   └── Processing.h             # организация основных вычислительных режимов
├── upper_bound_calculator/
│   ├── FullEnumUpperBoundCalculator.h # многопоточный перебор вершин
│   ├── LebesgueFunction.h       # вычисление значения функции Лебега в точке
│   └── UpperBoundCalculatorInterface.h # интерфейс для вычисления верхней границы
└── vertices_provider/
    ├── HadamardFileProvider.h   # загрузка матриц Адамара из файла
    ├── ISimplexProvider.h       # интерфейс для получения вершин симплекса
    ├── SilvesterMethodProvider.h # построение симплекса методом Сильвестра
    └── SimplexProviderFactory.h # фабрика для создания провайдеров

src/
├── basis/
│   └── SimplexBasis.cpp
├── cli/
│   └── CommandLineParser.cpp
├── hadamard_utils/
│   ├── HadamardMatrixIterator.cpp
│   └── HadamardUtils.cpp
├── processing/
│   └── Processing.cpp
├── upper_bound_calculator/
│   ├── FullEnumUpperBoundCalculator.cpp
│   └── LebesgueFunction.cpp
└── vertices_provider/
    ├── HadamardFileProvider.cpp
    ├── SilvesterMethodProvider.cpp
    └── SimplexProviderFactory.cpp
\end{verbatim}

\subsection*{Компиляция и запуск}

Для компиляции программы необходимо установить библиотеку Eigen и систему сборки CMake. Команды для сборки:

\begin{verbatim}
mkdir build && cd build
cmake -DCMAKE_BUILD_TYPE=Release ..
make -j$(nproc)
\end{verbatim}

Примеры запуска:

\begin{verbatim}
# Стандартное построение (Сильвестр) для размерности 31
./HadamardSimplexNorm --dimension 31

# Пакетная обработка всех файлов в каталоге
./HadamardSimplexNorm --batch ../input --output results.txt

# Загрузка матрицы Адамара из файла
./HadamardSimplexNorm --load-hadamard input/hadamard_32.txt --dimension 31
\end{verbatim}

\newpage

\section*{Приложение B. Заголовочные файлы}
\addcontentsline{toc}{section}{Приложение B. Заголовочные файлы}

\begin{lstlisting}
// Файл: ./src/main.cpp

#include "cli/CommandLineParser.h"
#include "../include/processing/Processing.h"
#include <iostream>

int main(int argc, char* argv[]) {
//  РАЗБОР АРГУМЕНТОВ КОМАНДНОЙ СТРОКИ
    ProgramOptions opts = CommandLineParser::parse(argc, argv);

    if (opts.showHelp) {
        CommandLineParser::printHelp(argv[0]);
        return 0;
    }

    if (!opts.partial) {
        std::cout << "Режим: полный перебор всех вершин" << std::endl;
    }

    if (opts.batchMode && !opts.inputDir.empty()) {
//  ПАКЕТНАЯ ОБРАБОТКА
        return Processing::processBatch(opts);
    }

    if (opts.loadFromFile && !opts.inputFile.empty()) {
//  ОБРАБОТКА ОДИНОЧНОГО ФАЙЛА
        return Processing::processSingleFile(opts.inputFile, opts);
    }

// СТАНДАРТНОЕ ПОСТРОЕНИЕ (СИЛЬВЕСТР)
    return Processing::processSilvester(opts);
}

=========================================================
// Файл: ./include/basis/SimplexBasis.h

#pragma once
#include <Eigen/Dense>
#include <vector>

class SimplexBasis {
   public:
    explicit SimplexBasis(const std::vector<Eigen::RowVectorXd>& vertices);
    [[nodiscard]] int    size() const;
    [[nodiscard]] double value(int i, const Eigen::RowVectorXd& x) const;
    static bool          verifyInterpolation(const SimplexBasis& basis, double tolerance);
    [[nodiscard]] const std::vector<Eigen::RowVectorXd>& getVertices() const {
        return vertices_;
    }

   private:
    std::vector<Eigen::RowVectorXd> vertices_;
    int                             dim_;
    Eigen::MatrixXd                 Ainv_;

   public:
    [[nodiscard]] Eigen::MatrixXd getAinv() const {
        return Ainv_;
    }
};

=========================================================
// Файл: ./include/cli/CommandLineParser.h

#pragma once
#include <string>
#include <cstdint>

struct ProgramOptions {
// Режим перебора
    bool    partial     = false;
    int64_t startVertex = 0;
    int64_t endVertex   = 0;
    int     dimension   = 7;

// Параметры выполнения
    int numThreads = 0;  // 0 = автоматически (hardware_concurrency)
    int64_t logInterval = 10'000'000;  // интервал логирования (количество вершин)

// Режим загрузки данных
    bool        loadFromFile = false;
    bool        batchMode    = false;
    std::string inputFile;
    std::string inputDir;
    enum class FileType { HADAMARD } fileType = FileType::HADAMARD;

// Помощь
    bool showHelp = false;
};

class CommandLineParser {
   public:
    static ProgramOptions parse(int argc, char* argv[]);
    static void           printHelp(const char* programName);

   private:
    static void parseRange(ProgramOptions& opts, int& i, char* argv[]);
    static void parseDimension(ProgramOptions& opts, int& i, char* argv[]);
    static void parseThreads(ProgramOptions& opts, int& i, char* argv[]);
    static void parseLogInterval(ProgramOptions& opts, int& i, char* argv[]);
};

=========================================================
// Файл: ./include/hadamard_utils/HadamardMatrixIterator.h

#pragma once
#include <string>
#include <vector>
#include <fstream>
#include <Eigen/Dense>

class HadamardMatrixIterator {
   public:
    explicit HadamardMatrixIterator(const std::string& filename);

    void            reset();
    bool            hasNext() const;
    Eigen::MatrixXd next();

   private:
    std::string   filename_;
    std::ifstream file_;
    int           currentOrder_;

    static bool isCommentLine(const std::string& line);
    static bool isSymbolicLine(const std::string& line);
    static bool isNumericLine(const std::string& line);

    static std::vector<double> parseLine(const std::string& line);
    Eigen::MatrixXd            readNextMatrix();
};

=========================================================
// Файл: ./include/hadamard_utils/HadamardUtils.h

#pragma once
#include <Eigen/Dense>
#include <vector>

class HadamardUtils {
   public:
    static std::vector<Eigen::RowVectorXd> getVerticesWithLogging(const Eigen::MatrixXd& H,
                                                                  int                    dim);
    static bool isHadamardMatrix(const Eigen::MatrixXd& H, double tolerance = 1e-10);

   private:
// Преобразование матрицы Адамара в (0/1)-матрицу D
    static Eigen::MatrixXd hadamardTo01Matrix(const Eigen::MatrixXd& H);

// Построение вершин симплекса из (0/1)-матрицы D
    static std::vector<Eigen::RowVectorXd> buildVerticesFrom01Matrix(const Eigen::MatrixXd& D);
};

=========================================================
// Файл: ./include/processing/Processing.h

#pragma once
#include <string>
#include <Eigen/Dense>

struct ProgramOptions;

class Processing {
    static int    processHadamardMatrixFromFile(const std::string&    filename,
                                                const ProgramOptions& opts, double& globalMin,
                                                int&                globalMinMatrixIdx,
                                                Eigen::RowVectorXd& globalMinPoint);
    static double processHadamardMatrix(const Eigen::MatrixXd& H, const ProgramOptions& opts,
                                        int matrixCounter);

   public:
    static int processBatch(const ProgramOptions& opts);
    static int processSingleFile(const std::string& filename, const ProgramOptions& opts);
    static int processSilvester(const ProgramOptions& opts);
};

=========================================================
// Файл: ./include/upper_bound_calculator/FullEnumUpperBoundCalculator.h

#pragma once
#include "upper_bound_calculator/UpperBoundCalculatorInterface.h"
#include "upper_bound_calculator/LebesgueFunction.h"
#include <atomic>
#include <Eigen/Dense>

class FullEnumUpperBoundCalculator : public UpperBoundCalculatorInterface {
   public:
    FullEnumUpperBoundCalculator(const LebesgueFunction& func, int dim);

    FullEnumUpperBoundCalculator(const LebesgueFunction& func, int dim, int64_t startVertex,
                                 int64_t endVertex);

// Запуск перебора. Возвращает максимальное значение суммы модулей базисных многочленов
    double calculate() override;

    [[nodiscard]] Eigen::RowVectorXd getMaxPoint() const override {
        return maxPoint_;
    }

   private:
    const LebesgueFunction& func_;
    int                     dim_;
    Eigen::RowVectorXd      maxPoint_;
    std::atomic<double>     globalMax_;
    std::atomic<int64_t>    processedVertices_;
    std::atomic<int64_t>    lastLogged_;
    std::atomic<int64_t>    currentBestVertex_;

// Параметры частичного перебора
    int64_t startVertex_;
    int64_t endVertex_;
    bool    partial_;
};

=========================================================
// Файл: ./include/upper_bound_calculator/LebesgueFunction.h

#pragma once
#include "basis/SimplexBasis.h"

class LebesgueFunction {
   public:
    explicit LebesgueFunction(const SimplexBasis& basis);
    double operator()(const Eigen::RowVectorXd& x) const;

   private:
    const SimplexBasis& basis_;
};

=========================================================
// Файл: ./include/upper_bound_calculator/UpperBoundCalculatorInterface.h

#pragma once
#include <Eigen/Dense>

class UpperBoundCalculatorInterface {
   public:
    virtual ~UpperBoundCalculatorInterface() = default;

// Запуск оптимизации. Возвращает максимальное значение функции
    virtual double calculate() = 0;

    void setNumThreads(int numThreads) {
        numThreads_ = numThreads;
    }
    void setLogInterval(int64_t interval) {
        logInterval_ = interval;
    }

// Возвращает точку куба, на которой достигнут максимум
    [[nodiscard]] virtual Eigen::RowVectorXd getMaxPoint() const = 0;

   protected:
    int     numThreads_  = 0;
    int64_t logInterval_ = 10'000'000;
};


=========================================================
// Файл: ./include/vertices_provider/HadamardFileProvider.h

#pragma once
#include "vertices_provider/ISimplexProvider.h"
#include "hadamard_utils/HadamardMatrixIterator.h"

#include <string>

class HadamardFileProvider : public ISimplexProvider {
   public:
    explicit HadamardFileProvider(const std::string& filename);
    std::vector<Eigen::RowVectorXd> getVertices() override;
    int                             getDimension() const override {
        return dim_;
    }
    HadamardMatrixIterator getIterator(const std::string& filename);

   private:
    static Eigen::MatrixXd readHadamardFile(const std::string& filename);
    Eigen::MatrixXd        H_;
    std::string            filename_;
    int                    dim_;
};

=========================================================
// Файл: ./include/vertices_provider/ISimplexProvider.h

#pragma once
#include <Eigen/Dense>
#include <vector>

class ISimplexProvider {
   public:
    virtual ~ISimplexProvider()                                  = default;
    virtual std::vector<Eigen::RowVectorXd> getVertices()        = 0;
    virtual int                             getDimension() const = 0;
};

=========================================================
// Файл: ./include/vertices_provider/SilvesterMethodProvider.h

#pragma once
#include "vertices_provider/ISimplexProvider.h"

class SilvesterMethodProvider : public ISimplexProvider {
   public:
    explicit SilvesterMethodProvider(int dim);
    std::vector<Eigen::RowVectorXd> getVertices() override;
    ;
    int getDimension() const override {
        return dim_;
    };
    static Eigen::MatrixXd hadamardMatrix(int order);

   private:
    int             dim_;
    Eigen::MatrixXd H_;
};

=========================================================
// Файл: ./include/vertices_provider/SimplexProviderFactory.h

#pragma once
#include "vertices_provider/ISimplexProvider.h"
#include <memory>
#include <string>

class SimplexProviderFactory {
   public:
    enum class Type { SILVESTER, FROM_HADAMARD_FILE };

    static std::unique_ptr<ISimplexProvider> create(Type type, int dim = 0,
                                                    const std::string& filename = "");};
\end{lstlisting}

\newpage

\section*{Приложение C. Реализация}
\addcontentsline{toc}{section}{Приложение C. Реализация}

\begin{lstlisting}
// Файл: ./src/basis/SimplexBasis.cpp

#include "basis/SimplexBasis.h"
#include <stdexcept>
#include <iostream>

SimplexBasis::SimplexBasis(const std::vector<Eigen::RowVectorXd>& vertices) : vertices_(vertices) {
    dim_ = (int)vertices.size() - 1;
    if ((int)vertices[0].size() != dim_)
        throw std::runtime_error("Несоответствие размерности вершин");

    Eigen::MatrixXd A(dim_ + 1, dim_ + 1);
    for (int i = 0; i <= dim_; ++i) {
        for (int j = 0; j < dim_; ++j) {
            A(i, j) = vertices[i][j];
        }
        A(i, dim_) = 1.0;
    }

    Ainv_ = A.inverse();

    std::cout << "A = \n" << A << "\n" << std::endl;
    std::cout << "Ainv = \n" << Ainv_ << "\n" << std::endl;
}

int SimplexBasis::size() const {
    return dim_ + 1;
}

double SimplexBasis::value(int i, const Eigen::RowVectorXd& x) const {
    if (x.size() != dim_)
        throw std::runtime_error("Несоответствие размерности точки");

    Eigen::RowVectorXd aug(dim_ + 1);
    for (int k = 0; k < dim_; ++k)
        aug(k) = x(k);
    aug(dim_) = 1.0;

    double res = 0.0;
    for (int k = 0; k <= dim_; ++k) {
        res += aug(k) * Ainv_(k, i);
    }
    return res;
}

bool SimplexBasis::verifyInterpolation(const SimplexBasis& basis, double tolerance) {
    const auto& vertices = basis.getVertices();
    double      maxError = 0.0;

    std::cout << "Проверка интерполяции: ℓ_i(v_j) должно быть 1 при i=j и 0 иначе" << std::endl;

    for (int i = 0; i < basis.size(); ++i) {
        for (int j = 0; j < basis.size(); ++j) {
            double val      = basis.value(i, vertices[j]);
            double expected = (i == j) ? 1.0 : 0.0;
            double error    = std::abs(val - expected);
            if (error > maxError)
                maxError = error;
            // Отладка: выводим первые несколько
            if (i < 3 && j < 3) {
                std::cout << "ℓ_" << i << "(v_" << j << ") = " << val;
                std::cout << " (ожидалось " << expected << ")" << std::endl;
            }
        }
    }

    std::cout << "Максимальная ошибка по всем парам (вершина, базис): " << maxError << std::endl;
    std::cout << std::endl;

    return maxError < tolerance;
}

=========================================================
// Файл: ./src/cli/CommandLineParser.cpp

#include "cli/CommandLineParser.h"
#include <iostream>
#include <string>

ProgramOptions CommandLineParser::parse(int argc, char* argv[]) {
    ProgramOptions opts;

    for (int i = 1; i < argc; ++i) {
        std::string arg = argv[i];

        if (arg == "-h" || arg == "--help") {
            opts.showHelp = true;
            return opts;
        }
        if (arg == "--dimension" && i + 1 < argc) {
            parseDimension(opts, i, argv);
            opts.endVertex = 1LL << opts.dimension;
        }
        if (arg == "--range" && i + 2 < argc) {
            parseRange(opts, i, argv);
        }
        if (arg == "--threads" && i + 1 < argc) {
            parseThreads(opts, i, argv);
        }
        if (arg == "--log-interval" && i + 1 < argc) {
            parseLogInterval(opts, i, argv);
        }
        if (arg == "--batch" && i + 1 < argc) {
            opts.batchMode = true;
            opts.inputDir  = argv[++i];
            std::cout << "Пакетный режим. Каталог: " << opts.inputDir << std::endl;
        }
        if (arg == "--load-hadamard" && i + 1 < argc) {
            opts.loadFromFile = true;
            opts.fileType     = ProgramOptions::FileType::HADAMARD;
            opts.inputFile    = argv[++i];
        }
    }

    return opts;
}

void CommandLineParser::parseRange(ProgramOptions& opts, int& i, char* argv[]) {
    opts.partial     = true;
    opts.startVertex = std::stoll(argv[++i]);
    opts.endVertex   = std::stoll(argv[++i]);
    std::cout << "Режим: частичный перебор" << std::endl;
    std::cout << "Диапазон вершин: [" << opts.startVertex << ", " << opts.endVertex << ")"
              << std::endl;
}

void CommandLineParser::parseDimension(ProgramOptions& opts, int& i, char* argv[]) {
    opts.dimension = std::stoi(argv[++i]);
    opts.endVertex = 1LL << opts.dimension;
    std::cout << "Размерность установлена: " << opts.dimension << std::endl;
    std::cout << "Всего вершин: " << opts.endVertex << " (2^" << opts.dimension << ")" << std::endl;
}

void CommandLineParser::parseThreads(ProgramOptions& opts, int& i, char* argv[]) {
    opts.numThreads = std::stoi(argv[++i]);
    std::cout << "Количество потоков: " << opts.numThreads << std::endl;
}

void CommandLineParser::parseLogInterval(ProgramOptions& opts, int& i, char* argv[]) {
    opts.logInterval = std::stoll(argv[++i]);
    std::cout << "Интервал логирования: " << opts.logInterval << " вершин" << std::endl;
}

void CommandLineParser::printHelp(const char* programName) {
    std::cout << "Использование: " << programName << " [ОПЦИИ]" << std::endl;
    std::cout << std::endl;
    std::cout << "ОПЦИИ РАЗМЕРНОСТИ:" << std::endl;
    std::cout << "  --dimension N       установить размерность пространства (по умолчанию 7)"
              << std::endl;
    std::cout << std::endl;
    std::cout << "ОПЦИИ ПЕРЕБОРА:" << std::endl;
    std::cout << "  --range START END   частичный перебор вершин в диапазоне [START, END)"
              << std::endl;
    std::cout << std::endl;
    std::cout << "ОПЦИИ ВЫПОЛНЕНИЯ:" << std::endl;
    std::cout << "  --threads N         количество потоков (0 = автоматически)" << std::endl;
    std::cout
        << "  --log-interval N    интервал логирования (количество вершин, по умолчанию 100000)"
        << std::endl;
    std::cout << std::endl;
    std::cout << "ОПЦИИ ЗАГРУЗКИ ДАННЫХ:" << std::endl;
    std::cout << "  --load-hadamard FILE  загрузить матрицу Адамара из файла" << std::endl;
    std::cout << "  --batch DIR           пакетная обработка всех файлов в каталоге" << std::endl;
    std::cout << std::endl;
    std::cout << "ПРИМЕРЫ:" << std::endl;
    std::cout << "  " << programName
              << "                                              # стандартный режим (dim=7)"
              << std::endl;
    std::cout << "  " << programName << " --dimension 31                               # dim=31"
              << std::endl;
    std::cout << "  " << programName << " --dimension 31 --threads 16                  # 16 потоков"
              << std::endl;
    std::cout << "  " << programName
              << " --dimension 31 --log-interval 1000000        # логирование каждые 1 млн вершин"
              << std::endl;
    std::cout << "  " << programName
              << " --range 0 500000000                          # частичный перебор" << std::endl;
    std::cout << "  " << programName
              << " --load-hadamard ../preset/test_Hadamar_8.txt # загрузка матрицы Адамара"
              << std::endl;
}

=========================================================
// Файл: ./src/hadamard_utils/HadamardMatrixIterator.cpp

#include "hadamard_utils/HadamardMatrixIterator.h"
#include <sstream>
#include <stdexcept>
#include <iostream>
#include <cctype>
#include <cmath>

HadamardMatrixIterator::HadamardMatrixIterator(const std::string& filename)
    : filename_(filename), currentOrder_(0) {
    reset();
}

void HadamardMatrixIterator::reset() {
    file_.close();
    file_.open(filename_);
    if (!file_.is_open()) {
        throw std::runtime_error("Не удалось открыть файл: " + filename_);
    }
}

bool HadamardMatrixIterator::hasNext() const {
    return file_.is_open() && !file_.eof();
}

bool HadamardMatrixIterator::isCommentLine(const std::string& line) {
    if (line.empty())
        return false;
    char c = line[0];
    // Строка комментария: первый символ не +, -, 0, 1, и не минус (который может быть началом -1)
    // Но -1 начинается с '-', поэтому нужно отличать числовой -1 от комментария
    // Проще: комментарий — если первый символ буква или не является частью допустимого формата
    if (c == '+' || c == '-')
        return false;
    if (c == '0' || c == '1')
        return false;
    if (c == ' ' || c == '\t')
        return false;
    return true;
}

bool HadamardMatrixIterator::isSymbolicLine(const std::string& line) {
    if (line.empty())
        return false;
    // Символьный формат: строка состоит только из '+' и '-'
    for (char ch : line) {
        if (ch != '+' && ch != '-')
            return false;
    }
    return true;
}

bool HadamardMatrixIterator::isNumericLine(const std::string& line) {
    if (line.empty())
        return false;
    // Числовой формат: строка содержит числа -1, 0 или 1, разделённые пробелами
    std::istringstream iss(line);
    double             val;
    while (iss >> val) {
        if (std::abs(val) > 1.01)
            return false;  // не -1, 0 или 1
    }
    return true;
}

std::vector<double> HadamardMatrixIterator::parseLine(const std::string& line) {
    std::vector<double> row;

    if (isSymbolicLine(line)) {
        for (char ch : line) {
            if (ch == '+')
                row.push_back(1.0);
            else if (ch == '-')
                row.push_back(-1.0);
        }
    } else {
        std::istringstream iss(line);
        double             val;
        while (iss >> val) {
            // Нормализуем к -1, 0, 1 (на случай -0.999999)
            if (std::abs(val - 1.0) < 1e-12)
                val = 1.0;
            else if (std::abs(val + 1.0) < 1e-12)
                val = -1.0;
            else if (std::abs(val) < 1e-12)
                val = 0.0;
            row.push_back(val);
        }
    }

    return row;
}

Eigen::MatrixXd HadamardMatrixIterator::readNextMatrix() {
    std::vector<std::vector<double>> rows;

    std::string line;
    while (std::getline(file_, line)) {
        // Удаляем пробелы
        size_t start = line.find_first_not_of(" \t\r\n");
        if (start == std::string::npos) {
            // Пустая строка — разделитель матриц, выходим
            break;
        }
        line = line.substr(start);

        size_t end = line.find_last_not_of(" \t\r\n");
        if (end != std::string::npos) {
            line = line.substr(0, end + 1);
        }

        if (line.empty())
            break;

        if (isCommentLine(line)) {
            continue;
        }

        // Строка данных матрицы
        if (isSymbolicLine(line) || isNumericLine(line)) {
            std::vector<double> row = parseLine(line);
            rows.push_back(row);
        } else {
            // Неизвестный формат — вероятно, разделитель
            break;
        }
    }

    int n = rows.size();
    if (n == 0) {
        throw std::runtime_error("Не удалось прочитать матрицу");
    }

    int m = rows[0].size();
    if (n != m) {
        throw std::runtime_error("Матрица не квадратная: " + std::to_string(n) + "x" +
                                 std::to_string(m));
    }

    Eigen::MatrixXd H(n, n);
    for (int i = 0; i < n; ++i) {
        for (int j = 0; j < n; ++j) {
            H(i, j) = rows[i][j];
        }
    }

    currentOrder_ = n;
    return H;
}

Eigen::MatrixXd HadamardMatrixIterator::next() {
    return readNextMatrix();
}

=========================================================
// Файл: ./src/hadamard_utils/HadamardUtils.cpp

#include "hadamard_utils/HadamardUtils.h"

#include <iostream>

Eigen::MatrixXd HadamardUtils::hadamardTo01Matrix(const Eigen::MatrixXd& H) {
    int             n = H.rows();
    Eigen::MatrixXd B = H;

    // Шаг 1: каждый столбец, начинающийся с -1, умножается на -1
    for (int j = 0; j < n; ++j) {
        if (B(0, j) == -1.0) {
            for (int i = 0; i < n; ++i) {
                B(i, j) = -B(i, j);
            }
        }
    }

    // Шаг 2: каждая строка, начинающаяся с -1, умножается на -1
    for (int i = 0; i < n; ++i) {
        if (B(i, 0) == -1.0) {
            for (int j = 0; j < n; ++j) {
                B(i, j) = -B(i, j);
            }
        }
    }

    // Шаг 3: подматрица C (строки и столбцы 2..n), замена: 1 → 0, -1 → 1
    int             dim = n - 1;
    Eigen::MatrixXd D(dim, dim);

    for (int i = 0; i < dim; ++i) {
        for (int j = 0; j < dim; ++j) {
            double val = B(i + 1, j + 1);
            D(i, j)    = (val == -1.0) ? 1.0 : 0.0;
        }
    }

    return D;
}

std::vector<Eigen::RowVectorXd> HadamardUtils::buildVerticesFrom01Matrix(const Eigen::MatrixXd& D) {
    int dim         = D.rows();
    int numVertices = dim + 1;

    std::vector<Eigen::RowVectorXd> vertices;
    vertices.reserve(numVertices);

    // Нулевая вершина
    Eigen::RowVectorXd zero(dim);
    zero.setZero();
    vertices.push_back(zero);

    // Строки матрицы D как вершины
    for (int i = 0; i < dim; ++i) {
        vertices.push_back(D.row(i));
    }

    return vertices;
}

std::vector<Eigen::RowVectorXd> HadamardUtils::getVerticesWithLogging(const Eigen::MatrixXd& H,
                                                                      int                    dim) {
    int order = dim + 1;

    std::cout << "=== Матрица Адамара H (1/-1): ===\n" << H << "\n" << std::endl;

    long double detH           = H.determinant();
    long double theoreticalDet = std::pow(order, order / 2);
    std::cout << "det(H) = " << detH << std::endl;
    std::cout << "Теоретическое значение: ±" << theoreticalDet << "\n" << std::endl;

    // Проверка определителя матрицы Адамара
    if (std::abs(std::abs(detH) - std::abs(theoreticalDet)) > 10) {
        std::cout << "ОШИБКА: Определитель матрицы Адамара (-1/1) не соответствует теоретическому "
                     "значению!"
                  << std::endl;
        std::cout << "Получено: " << detH << ", ожидается ±" << theoreticalDet << std::endl;
        std::exit(1);
    }

    Eigen::MatrixXd D = hadamardTo01Matrix(H);

    std::cout << "=== Матрица D (0/1): ===\n" << D << "\n" << std::endl;

    double detD        = D.determinant();
    double expectedDet = std::abs(detH) / std::pow(2.0, dim);
    std::cout << "det(D) = " << detD << std::endl;
    std::cout << "Теоретическое значение: ±" << expectedDet << "\n" << std::endl;

    // Проверка определителя (0/1)-матрицы
    if (std::abs(std::abs(detD) - std::abs(expectedDet)) > 10) {
        std::cout << "ОШИБКА: Определитель матрицы Адамара (0/1) не соответствует теоретическому "
                     "значению!"
                  << std::endl;
        std::cout << "Получено: " << detD << ", ожидается ±" << expectedDet << std::endl;
        std::exit(1);
    }

    auto vertices = buildVerticesFrom01Matrix(D);

    std::cout << "Вершины симплекса:" << std::endl;
    for (size_t i = 0; i < vertices.size(); ++i) {
        std::cout << "v[" << i << "]: " << vertices[i] << std::endl;
    }

    std::cout << std::endl;

    return vertices;
}

bool HadamardUtils::isHadamardMatrix(const Eigen::MatrixXd& H, double tolerance) {
    int n = H.rows();

    // Шаг 1: проверка квадратности
    if (H.cols() != n) {
        std::cerr << "Ошибка: матрица не квадратная" << std::endl;
        return false;
    }

    // Шаг 2: проверка элементов (±1)
    for (int i = 0; i < n; ++i) {
        for (int j = 0; j < n; ++j) {
            double val = std::abs(H(i, j));
            if (std::abs(val - 1.0) > tolerance) {
                std::cerr << "Ошибка: элемент [" << i << "][" << j << "] = " << H(i, j)
                          << " не равен ±1" << std::endl;
                return false;
            }
        }
    }

    // Шаг 3: проверка ортогональности строк (H * H^T = n * I)
    Eigen::MatrixXd HHT = H * H.transpose();
    for (int i = 0; i < n; ++i) {
        for (int j = 0; j < n; ++j) {
            double expected = (i == j) ? n : 0.0;
            if (std::abs(HHT(i, j) - expected) > tolerance) {
                std::cerr << "Ошибка: (H*H^T)[" << i << "][" << j << "] = " << HHT(i, j)
                          << ", ожидается " << expected << std::endl;
                return false;
            }
        }
    }

    return true;
}
=========================================================
// Файл: ./src/processing/Processing.cpp

#include "processing/Processing.h"

#include "basis/SimplexBasis.h"
#include "cli/CommandLineParser.h"
#include "vertices_provider/SimplexProviderFactory.h"
#include "upper_bound_calculator/FullEnumUpperBoundCalculator.h"
#include "upper_bound_calculator/LebesgueFunction.h"
#include "hadamard_utils/HadamardMatrixIterator.h"
#include "hadamard_utils/HadamardUtils.h"

#include <chrono>
#include <iomanip>
#include <iostream>
#include <filesystem>

namespace fs = std::filesystem;

struct ProgramOptions;

// Функция обработки одной матрицы Адамара
double Processing::processHadamardMatrix(const Eigen::MatrixXd& H, const ProgramOptions& opts,
                                         int matrixCounter) {
    int order         = H.rows();
    int expectedOrder = opts.dimension + 1;

    // Проверка порядка матрицы
    if (order != expectedOrder) {
        std::cout << "Предупреждение: порядок матрицы " << order << " не соответствует ожидаемому "
                  << expectedOrder << std::endl;
        std::cout << "Матрица #" << matrixCounter << ": неверный порядок (" << order << " вместо "
                  << expectedOrder << ")" << std::endl;
        return -1.0;
    }

    std::cout << "\n" << std::string(60, '=') << std::endl;
    std::cout << "Обработка матрицы #" << matrixCounter << " (порядок " << order << ")"
              << std::endl;
    std::cout << std::string(60, '=') << std::endl;

    // Строим симплекс из матрицы Адамара
    auto vertices = HadamardUtils::getVerticesWithLogging(H, order - 1);

    SimplexBasis     basis(vertices);
    LebesgueFunction lebesgueFunc(basis);

    // Проверка интерполяции
    if (!SimplexBasis::verifyInterpolation(basis, 0.00001)) {
        std::cout << "ОШИБКА: Интерполяционные свойства не выполняются!" << std::endl;
        std::cout << "Матрица #" << matrixCounter << ": ошибка интерполяции" << std::endl;
        return -1.0;
    }

    // Вычисление верхней границы
    FullEnumUpperBoundCalculator calculator(lebesgueFunc, opts.dimension);
    calculator.setNumThreads(opts.numThreads);
    calculator.setLogInterval(opts.logInterval);

    auto   start       = std::chrono::high_resolution_clock::now();
    double maxLebesgue = calculator.calculate();
    auto   end         = std::chrono::high_resolution_clock::now();
    auto   duration    = std::chrono::duration_cast<std::chrono::seconds>(end - start);

    Eigen::RowVectorXd maxPoint = calculator.getMaxPoint();

    std::cout << "Максимальная сумма модулей: " << std::fixed << std::setprecision(6)
              << maxLebesgue << std::endl;
    std::cout << "Время выполнения: " << duration.count() << " секунд" << std::endl;
    std::cout << std::endl;

    // Вывод результата
    std::cout << "ИТОГ для Матрица #" << matrixCounter << ", порядок " << order << ": "
              << std::fixed << std::setprecision(6) << maxLebesgue << " (за " << duration.count()
              << " сек)" << std::endl;

    return maxLebesgue;
}

// Функция обработки одного файла (может содержать несколько матриц)
int Processing::processHadamardMatrixFromFile(const std::string&    filename,
                                              const ProgramOptions& opts, double& globalMin,
                                              int&                globalMinMatrixIdx,
                                              Eigen::RowVectorXd& globalMinPoint) {
    int                matrixCounter = 1;
    double             fileMin       = std::numeric_limits<double>::max();
    int                fileMinIdx    = -1;
    Eigen::RowVectorXd fileMinPoint;

    std::cout << "\n" << std::string(80, '=') << std::endl;
    std::cout << "Обработка файла: " << filename << std::endl;
    std::cout << std::string(80, '=') << std::endl;

    try {
        HadamardMatrixIterator it(filename);
        while (it.hasNext()) {
            Eigen::MatrixXd H        = it.next();
            double          lebesgue = processHadamardMatrix(H, opts, matrixCounter);

            if (lebesgue >= 0) {
                if (lebesgue < fileMin) {
                    fileMin    = lebesgue;
                    fileMinIdx = matrixCounter;
                }
            }
            matrixCounter++;
        }
    } catch (const std::exception& e) {
        std::cout << "Ошибка при обработке файла " << filename << ": " << e.what() << std::endl;
        std::cout << "Ошибка в файле " << filename << ": " << e.what() << std::endl;
    }

    if (fileMinIdx != -1) {
        std::cout << "\n>>> МИНИМУМ в файле " << filename << ": матрица #" << fileMinIdx
                  << " дала значение " << std::fixed << std::setprecision(4) << fileMin
                  << std::endl;

        // Обновляем глобальный минимум
        if (fileMin < globalMin) {
            globalMin          = fileMin;
            globalMinMatrixIdx = fileMinIdx;
            globalMinPoint     = fileMinPoint;
        }
    }

    return 0;
}

int Processing::processBatch(const ProgramOptions& opts) {
    // Собираем все .txt файлы в каталоге
    std::vector<std::string> files;
    if (fs::exists(opts.inputDir) && fs::is_directory(opts.inputDir)) {
        for (const auto& entry : fs::directory_iterator(opts.inputDir)) {
            if (entry.is_regular_file() && entry.path().extension() == ".txt") {
                files.push_back(entry.path().string());
            }
        }
    } else {
        std::cerr << "Каталог не найден: " << opts.inputDir << std::endl;
        return 1;
    }

    std::sort(files.begin(), files.end());
    std::cout << "Найдено файлов для обработки: " << files.size() << std::endl;

    // Глобальный минимум
    double             globalMin          = std::numeric_limits<double>::max();
    int                globalMinMatrixIdx = -1;
    Eigen::RowVectorXd globalMinPoint;

    for (const auto& file : files) {
        processHadamardMatrixFromFile(file, opts, globalMin, globalMinMatrixIdx, globalMinPoint);
    }

    // Вывод глобального минимума
    if (globalMinMatrixIdx != -1) {
        std::cout << "\n" << std::string(80, '*') << std::endl;
        std::cout << "ГЛОБАЛЬНЫЙ МИНИМУМ: матрица #" << globalMinMatrixIdx << " дала значение "
                  << std::fixed << std::setprecision(6) << globalMin << std::endl;
        std::cout << std::string(80, '*') << std::endl;
    }

    return 0;
}

int Processing::processSingleFile(const std::string& filename, const ProgramOptions& opts) {
    double             dummyMin;
    int                dummyIdx;
    Eigen::RowVectorXd dummyPoint;
    return processHadamardMatrixFromFile(filename, opts, dummyMin, dummyIdx, dummyPoint);
}

int Processing::processSilvester(const ProgramOptions& opts) {
    auto vertices_provider =
        SimplexProviderFactory::create(SimplexProviderFactory::Type::SILVESTER, opts.dimension, "");
    auto vertices = vertices_provider->getVertices();

    std::cout << "Число вершин симплекса: " << vertices.size() << "\n" << std::endl;

    SimplexBasis     basis(vertices);
    LebesgueFunction lebesgueFunc(basis);

    std::cout << "\n=== Проверка интерполяции ===" << std::endl;
    if (!SimplexBasis::verifyInterpolation(basis, 0.00001)) {
        std::cout << "ВНИМАНИЕ: Интерполяционные свойства не выполняются!" << std::endl;
        return 1;
    }
    std::cout << "ВЕРНО: Интерполяционные свойства выполняются" << "\n" << std::endl;

    FullEnumUpperBoundCalculator calculator(lebesgueFunc, opts.dimension);
    calculator.setNumThreads(opts.numThreads);
    calculator.setLogInterval(opts.logInterval);

    auto   start       = std::chrono::high_resolution_clock::now();
    double maxLebesgue = calculator.calculate();
    auto   end         = std::chrono::high_resolution_clock::now();
    auto   duration    = std::chrono::duration_cast<std::chrono::seconds>(end - start);

    std::cout << "\n=== РЕЗУЛЬТАТЫ ===" << std::endl;
    std::cout << std::fixed << std::setprecision(6);
    std::cout << "Максимальная сумма модулей: " << maxLebesgue << std::endl;

    Eigen::RowVectorXd maxPoint = calculator.getMaxPoint();
    std::cout << "Найдена в вершине: [";
    for (int i = 0; i < maxPoint.size(); ++i) {
        std::cout << static_cast<int>(maxPoint(i) + 0.5);
        if (i + 1 < maxPoint.size())
            std::cout << " ";
    }
    std::cout << "]" << std::endl;

    std::cout << "Время выполнения: " << duration.count() << " секунд" << std::endl;

    std::cout << "Сильвестр (dim=" << opts.dimension << "): " << std::fixed << std::setprecision(6)
              << maxLebesgue << " (за " << duration.count() << " сек)" << std::endl;
    return 0;
}


=========================================================
// Файл: ./src/upper_bound_calculator/FullEnumUpperBoundCalculator.cpp

#include "upper_bound_calculator/FullEnumUpperBoundCalculator.h"
#include <thread>
#include <vector>
#include <iostream>
#include <iomanip>
#include <chrono>
#include <mutex>

static std::mutex cout_mutex;

FullEnumUpperBoundCalculator::FullEnumUpperBoundCalculator(const LebesgueFunction& func, int dim)
    : func_(func),
      dim_(dim),
      globalMax_(0.0),
      processedVertices_(0),
      lastLogged_(0),
      currentBestVertex_(0),
      startVertex_(0),
      endVertex_(0),
      partial_(false) {
    maxPoint_ = Eigen::RowVectorXd::Zero(dim_);
}

FullEnumUpperBoundCalculator::FullEnumUpperBoundCalculator(const LebesgueFunction& func, int dim,
                                                           int64_t startVertex, int64_t endVertex)
    : func_(func),
      dim_(dim),
      globalMax_(0.0),
      processedVertices_(0),
      lastLogged_(0),
      currentBestVertex_(0),
      startVertex_(startVertex),
      endVertex_(endVertex),
      partial_(true) {
    maxPoint_ = Eigen::RowVectorXd::Zero(dim_);
}

static void enumerateSubset(const LebesgueFunction* func, int dim, int64_t start, int64_t end,
                            std::atomic<double>* globalMax, std::atomic<int64_t>* processed,
                            std::atomic<int64_t>* lastLogged, std::atomic<int64_t>* bestVertex,
                            int64_t totalVertices, int64_t logInterval,
                            std::chrono::steady_clock::time_point startTime) {
    double  localMax        = 0.0;
    int64_t localBestVertex = 0;

    for (int64_t idx = start; idx < end; ++idx) {
        Eigen::RowVectorXd x(dim);
        for (int i = 0; i < dim; ++i) {
            x(i) = ((idx >> i) & 1) ? 1.0 : 0.0;
        }

        double val = (*func)(x);

        // Обновляем локальный и глобальный максимум сразу при нахождении лучшей вершины
        if (val > localMax) {
            localMax        = val;
            localBestVertex = idx;

            double currentMax = globalMax->load();
            while (localMax > currentMax) {
                if (globalMax->compare_exchange_weak(currentMax, localMax)) {
                    bestVertex->store(localBestVertex);
                    break;
                }
                currentMax = globalMax->load();  // Перечитываем при неудаче из-за другого потока
            }
        }

        int64_t processedCount = processed->fetch_add(1) + 1;

        // Логирование каждые logInterval итераций
        if (processedCount - lastLogged->load() >= logInterval) {
            lastLogged->store(processedCount);

            double percent =
                100.0 * static_cast<double>(processedCount) / static_cast<double>(totalVertices);
            auto now = std::chrono::steady_clock::now();
            auto elapsed =
                std::chrono::duration_cast<std::chrono::seconds>(now - startTime).count();
            auto elapsed_us =
                std::chrono::duration_cast<std::chrono::microseconds>(now - startTime).count();

            double us_per_iter =
                (processedCount > 0) ? static_cast<double>(elapsed_us) / processedCount : 0.0;
            double vertices_per_sec =
                (elapsed > 0) ? static_cast<double>(processedCount) / elapsed : 0.0;

            double remaining_sec = 0.0;
            if (processedCount > 0) {
                remaining_sec = static_cast<double>(elapsed) * (totalVertices - processedCount) /
                                processedCount;
            }

            auto time_t = std::chrono::system_clock::to_time_t(std::chrono::system_clock::now());

            std::lock_guard<std::mutex> lock(cout_mutex);
            std::cout << std::put_time(std::localtime(&time_t), "%H:%M:%S")
                      << " | Прогресс: " << std::fixed << std::setprecision(4) << percent << "%"
                      << " | Обработано: " << processedCount << " / " << totalVertices
                      << " | Прошло: " << elapsed << " сек"
                      << " | Скорость: " << std::setprecision(0) << vertices_per_sec << " в/сек"
                      << " | мкс/итер: " << std::setprecision(2) << us_per_iter
                      << " | Осталось ≈: " << static_cast<int>(remaining_sec) << " сек"
                      << " | Текущее: " << val << " | Рекорд: " << globalMax->load() << std::endl;
        }
    }
}

double FullEnumUpperBoundCalculator::calculate() {
    int64_t totalVertices = partial_ ? (endVertex_ - startVertex_) : (1LL << dim_);
    int64_t startOffset   = partial_ ? startVertex_ : 0;

    int numThreads = numThreads_;
    if (numThreads <= 0) {
        numThreads = std::thread::hardware_concurrency();
        if (numThreads == 0)
            numThreads = 8;
    }

    if (!partial_) {
        std::cout << "=== Полный перебор вершин гиперкуба ===" << std::endl;
        std::cout << "Размерность: " << dim_ << std::endl;
        std::cout << "Всего вершин: " << totalVertices << " (2^" << dim_ << ")" << std::endl;
    } else {
        std::cout << "=== Частичный перебор вершин гиперкуба ===" << std::endl;
        std::cout << "Размерность: " << dim_ << std::endl;
        std::cout << "Диапазон: [" << startVertex_ << ", " << endVertex_ << ")" << std::endl;
        std::cout << "Всего вершин: " << totalVertices << std::endl;
    }
    std::cout << "Потоков: " << numThreads << std::endl;
    std::cout << "Логирование каждые " << logInterval_ << " вершин" << std::endl;
    std::cout << "=========================================" << std::endl;

    auto startTime = std::chrono::steady_clock::now();

    int64_t                  chunkSize = totalVertices / numThreads;
    std::vector<std::thread> threads;

    for (int t = 0; t < numThreads; ++t) {
        int64_t start = startOffset + t * chunkSize;
        int64_t end   = (t == numThreads - 1) ? startOffset + totalVertices : start + chunkSize;

        threads.emplace_back(enumerateSubset, &func_, dim_, start, end, &globalMax_,
                             &processedVertices_, &lastLogged_, &currentBestVertex_, totalVertices,
                             logInterval_, startTime);
    }

    for (auto& th : threads) {
        th.join();
    }

    auto endTime  = std::chrono::steady_clock::now();
    auto duration = std::chrono::duration_cast<std::chrono::seconds>(endTime - startTime);

    std::cout << "\n=== Перебор завершён ===" << std::endl;
    std::cout << "Время: " << duration.count() << " секунд" << std::endl;
    std::cout << "Обработано вершин: " << processedVertices_.load() << std::endl;

    double avgSpeed = 0.0;
    if (duration.count() > 0) {
        avgSpeed = static_cast<double>(processedVertices_.load()) / duration.count();
    }
    std::cout << "Средняя скорость: " << std::fixed << std::setprecision(0) << avgSpeed << " в/сек"
              << std::endl;

    std::cout << std::fixed << std::setprecision(3);
    std::cout << "Максимальная сумма модулей: " << globalMax_.load() << std::endl;

    int64_t bestIdx = currentBestVertex_.load();
    for (int i = 0; i < dim_; ++i) {
        maxPoint_(i) = ((bestIdx >> i) & 1) ? 1.0 : 0.0;
    }

    return globalMax_.load();
}


=========================================================
// Файл: ./src/upper_bound_calculator/LebesgueFunction.cpp

#include "upper_bound_calculator/LebesgueFunction.h"
#include <cmath>

LebesgueFunction::LebesgueFunction(const SimplexBasis& basis) : basis_(basis) {
}

double LebesgueFunction::operator()(const Eigen::RowVectorXd& x) const {
    double sum = 0.0;
    for (int i = 0; i < basis_.size(); ++i) {
        sum += std::abs(basis_.value(i, x));
    }
    return sum;
}

=========================================================
// Файл: ./src/vertices_provider/HadamardFileProvider.cpp

#include "vertices_provider/HadamardFileProvider.h"

#include "hadamard_utils/HadamardMatrixIterator.h"
#include "hadamard_utils/HadamardUtils.h"
#include <fstream>
#include <iostream>
#include <sstream>
#include <stdexcept>

HadamardFileProvider::HadamardFileProvider(const std::string& filename)
    : filename_(filename), dim_(0) {
    H_ = readHadamardFile(filename_);

    int order = H_.rows();
    dim_      = order - 1;

    // Проверка, что загруженная матрица является матрицей Адамара
    if (!HadamardUtils::isHadamardMatrix(H_)) {
        throw std::runtime_error("Загруженная матрица не является матрицей Адамара!");
    }

    std::cout << "Матрица Адамара порядка " << order << " успешно загружена и проверена"
              << std::endl;
}

std::vector<Eigen::RowVectorXd> HadamardFileProvider::getVertices() {
    return HadamardUtils::getVerticesWithLogging(H_, dim_);
}

Eigen::MatrixXd HadamardFileProvider::readHadamardFile(const std::string& filename) {
    std::ifstream file(filename);
    if (!file.is_open()) {
        throw std::runtime_error("Не удалось открыть файл: " + filename);
    }

    std::vector<std::vector<double>> rows;
    std::string                      line;
    bool                             isSymbolicFormat = false;

    while (std::getline(file, line)) {
        // Удаляем пробелы в начале и конце
        line.erase(0, line.find_first_not_of(" \t\r\n"));
        line.erase(line.find_last_not_of(" \t\r\n") + 1);

        // Пропускаем пустые строки
        if (line.empty())
            continue;

        // Определяем формат по первому непустому символу
        char firstChar = line[0];

        if (firstChar == '+' || firstChar == '-') {
            isSymbolicFormat = true;

            // Парсим символьный формат
            std::vector<double> row;
            for (char ch : line) {
                if (ch == '+') {
                    row.push_back(1.0);
                } else if (ch == '-') {
                    row.push_back(-1.0);
                }
                // Игнорируем пробелы и другие символы
            }
            if (!row.empty()) {
                rows.push_back(row);
            }
        } else if (firstChar == '1' || (firstChar >= '0' && firstChar <= '9')) {
            isSymbolicFormat = false;

            std::istringstream  iss(line);
            std::vector<double> row;
            double              val;
            while (iss >> val) {
                row.push_back(val);
            }
            if (!row.empty()) {
                rows.push_back(row);
            }
        } else {
            // Строка-комментарий (начинается с буквы или другого символа)
            std::cout << "Пропущена строка-комментарий: " << line << std::endl;
        }
    }

    if (rows.empty()) {
        throw std::runtime_error("Файл не содержит данных");
    }

    int n = rows.size();
    int m = rows[0].size();

    if (n != m) {
        throw std::runtime_error("Матрица не квадратная: " + std::to_string(n) + "x" +
                                 std::to_string(m));
    }

    Eigen::MatrixXd H(n, n);
    for (int i = 0; i < n; ++i) {
        for (int j = 0; j < n; ++j) {
            H(i, j) = rows[i][j];
        }
    }

    std::cout << "Файл прочитан в " << (isSymbolicFormat ? "символьном" : "числовом") << " формате"
              << std::endl;

    return H;
}

HadamardMatrixIterator HadamardFileProvider::getIterator(const std::string& filename) {
    return HadamardMatrixIterator(filename);
}

=========================================================
// Файл: ./src/vertices_provider/SilvesterMethodProvider.cpp

#include "vertices_provider/SilvesterMethodProvider.h"
#include "hadamard_utils/HadamardUtils.h"
#include <iostream>
#include <stdexcept>
#include <cmath>

SilvesterMethodProvider::SilvesterMethodProvider(int dim) : dim_(dim) {
    int order = dim + 1;
    H_        = hadamardMatrix(order);

    if (!HadamardUtils::isHadamardMatrix(H_)) {
        throw std::runtime_error("Созданная алгоритмом матрица не является матрицей Адамара!");
    }
}

Eigen::MatrixXd SilvesterMethodProvider::hadamardMatrix(int order) {
    if (order == 1) {
        Eigen::MatrixXd H(1, 1);
        H(0, 0) = 1.0;
        return H;
    }
    if (order <= 0 || (order & (order - 1)) != 0) {
        throw std::runtime_error("Метод Сильвестра требует порядок, являющийся степенью двойки");
    }
    int             half  = order / 2;
    Eigen::MatrixXd Hhalf = hadamardMatrix(half);
    Eigen::MatrixXd H(order, order);
    H.topLeftCorner(half, half)     = Hhalf;
    H.topRightCorner(half, half)    = Hhalf;
    H.bottomLeftCorner(half, half)  = Hhalf;
    H.bottomRightCorner(half, half) = -Hhalf;
    return H;
}

std::vector<Eigen::RowVectorXd> SilvesterMethodProvider::getVertices() {
    return HadamardUtils::getVerticesWithLogging(H_, dim_);
}

=========================================================
// Файл: ./src/vertices_provider/SimplexProviderFactory.cpp

#include "vertices_provider/SimplexProviderFactory.h"
#include "vertices_provider/SilvesterMethodProvider.h"
#include "vertices_provider/HadamardFileProvider.h"

std::unique_ptr<ISimplexProvider> SimplexProviderFactory::create(Type type, int dim,
                                                                 const std::string& filename) {
    switch (type) {
        case Type::SILVESTER:
            return std::make_unique<SilvesterMethodProvider>(dim);
        case Type::FROM_HADAMARD_FILE:
            return std::make_unique<HadamardFileProvider>(filename);
        default:
            throw std::runtime_error("Неизвестный тип провайдера");
    }
}

\end{lstlisting}
\end{document}